% Part: incompleteness
% Chapter: theories-computability
% Section: introduction

\documentclass[../../../include/open-logic-section]{subfiles}

\begin{document}

\olfileid{inc}{tcp}{int}

\olsection{పరిచయం}

\begin{editorial}
ఈ విభాగాన్ని తిరిగి రాయాలి.
\end{editorial}

మనకు ఇప్పటికే కిందివి ఉన్నాయి:
\begin{enumerate}
\item ఒక ప్రమేయం~$\Th{Q}$లో ప్రాతినిధ్యం చేయదగినది అనడం అంటే ఏమిటో తెలిపే
  నిర్వచనం (\olref[req][int]{defn:representable-fn})
\item ఒక సంబంధం~$\Th{Q}$లో ప్రాతినిధ్యం చేయదగినది అనడం అంటే ఏమిటో తెలిపే
  నిర్వచనం (\olref[req][rel]{defn:representing-relations})
\item $\Th{Q}$లో ప్రాతినిధ్యం చేయదగిన ప్రమేయాలు సరిగ్గా గణనీయ ప్రమేయాలేనని
  చెప్పే సిద్ధాంతం (\olref[req][int]{thm:representable-iff-comp})
\item $\Th{Q}$లో ప్రాతినిధ్యం చేయదగిన సంబంధాలు సరిగ్గా గణనీయ సంబంధాలేనని
  చెప్పే సిద్ధాంతం (\olref[req][rel]{thm:representing-rels})
\end{enumerate}

ఒక {\em సిద్ధాంతం} అంటే నిగమన పరంగా సంవృతమైన వాక్యాల సమితి; అంటే $T$
$!A$ను నిరూపించినప్పుడల్లా $!A$, $T$లో ఉండాలి. కొన్ని వాక్యాలతో పాటు అవి
సూచించే విషయాలన్నిటినీ కలిపిన సమాహారంగా సిద్ధాంతాన్ని భావించడం బహుశా
ఉత్తమం. ఇక నుంచి \olref[req][int]{sec}లోని ఎనిమిది స్వీకృతాల నుంచి
వ్యుత్పాదించదగిన వాక్యాల సమితితో ఏర్పడే {\em సిద్ధాంతాన్ని} సూచించడానికి
$\Th{Q}$ను వాడతాం. $\Th{Q}$ సూత్రాలను సంఖ్యలుగా సంకేతీకరించవచ్చని
గుర్తుంచుకోండి; $!A$ అటువంటి సూత్రమైతే, $!A$ను సంకేతీకరించే సంఖ్యను
$\Gn{!A}$తో సూచిద్దాం. ఈ సంకేతీకరణను ఆధారంగా తీసుకుని, వివిధ సూత్రాల
సమితులు గణనీయాలా కాదా అని ఇప్పుడు అడగవచ్చు.

\end{document}
